% 硕士学位论文 LaTeX 模板
% 编译方式：xelatex -> bibtex -> xelatex -> xelatex
\documentclass[12pt,a4paper,openright,UTF8]{ctexbook}

% ==================== 基础宏包 ====================
\usepackage{geometry}
\geometry{left=2.5cm,right=2.5cm,top=2.5cm,bottom=2.5cm}

\usepackage{xeCJK}
\usepackage{amsmath,amssymb,amsfonts}
\usepackage{graphicx}
\usepackage{booktabs}
\usepackage{longtable}
\usepackage{multirow}
\usepackage{array}
\usepackage{float}
\usepackage{caption}
\usepackage{subcaption}
\usepackage{enumitem}
\usepackage{algorithm}
\usepackage{algpseudocode}
\usepackage{listings}
\usepackage{xcolor}
\usepackage{hyperref}
\usepackage{cleveref}
\usepackage{setspace}
\usepackage{fancyhdr}
\usepackage{titlesec}
\usepackage{tocloft}

% ==================== 参考文献 ====================
\usepackage[numbers,sort&compress]{natbib}
\bibliographystyle{gbt7714-numerical}

% ==================== 页面设置 ====================
\setlength{\headheight}{14pt}
\addtolength{\topmargin}{-2pt}
\pagestyle{fancy}
\fancyhf{}
\fancyhead[C]{\small \leftmark}
\fancyfoot[C]{\thepage}
\renewcommand{\headrulewidth}{0.4pt}

% 章节标题格式
\ctexset{
    chapter = {
        format = \zihao{-3}\bfseries\centering,
        beforeskip = 20pt,
        afterskip = 20pt
    },
    section = {
        format = \zihao{4}\bfseries,
        beforeskip = 12pt,
        afterskip = 6pt
    },
    subsection = {
        format = \zihao{-4}\bfseries,
        beforeskip = 8pt,
        afterskip = 4pt
    }
}

% 行距
\onehalfspacing

% 代码高亮
\lstset{
    language=Python,
    basicstyle=\small\ttfamily,
    keywordstyle=\color{blue},
    commentstyle=\color{gray},
    stringstyle=\color{orange},
    numbers=left,
    numberstyle=\tiny\color{gray},
    stepnumber=1,
    frame=single,
    breaklines=true,
    showstringspaces=false
}

% 图片路径
\graphicspath{{figures/}}

% ==================== 论文信息 ====================
\newcommand{\thesistitle}{论文题目}
\newcommand{\thesisetitle}{English Title of Your Thesis}
\newcommand{\authorname}{你的名字}
\newcommand{\studentid}{2024XXXXXX}
\newcommand{\majorname}{你的专业}
\newcommand{\supervisor}{导师姓名}
\newcommand{\institutename}{你的学院}
\newcommand{\universityname}{你的大学}
\newcommand{\defensedate}{2026年X月}

\begin{document}
\hypersetup{pageanchor=false}

% ==================== 封面 ====================
\begin{titlepage}
    \centering
    \vspace*{2cm}

    {\zihao{2}\bfseries \universityname}

    \vspace{0.8cm}
    {\zihao{-1}\bfseries 硕士学位论文}

    \vspace{2.5cm}
    {\zihao{2}\bfseries \thesistitle}

    \vspace{0.8cm}
    {\zihao{3} \thesisetitle}

    \vspace{3cm}
    \begin{tabular}{rl}
        \zihao{4}\textbf{作\quad 者：} & \zihao{4}\authorname \\[0.5cm]
        \zihao{4}\textbf{学\quad 号：} & \zihao{4}\studentid \\[0.5cm]
        \zihao{4}\textbf{专\quad 业：} & \zihao{4}\majorname \\[0.5cm]
        \zihao{4}\textbf{导\quad 师：} & \zihao{4}\supervisor \\[0.5cm]
        \zihao{4}\textbf{学\quad 院：} & \zihao{4}\institutename \\
    \end{tabular}

    \vfill
    {\zihao{4} \defensedate}
\end{titlepage}

% ==================== 原创性声明（可选） ====================
\newpage
\thispagestyle{empty}
\begin{center}
    {\zihao{3}\bfseries 学位论文原创性声明}
\end{center}

\vspace{1cm}
\zihao{-4}
本人郑重声明：所呈交的学位论文，是本人在导师的指导下，独立进行研究工作所取得的成果。除文中已经注明引用的内容外，本论文不包含任何其他个人或集体已经发表或撰写过的作品成果。对本文的研究做出重要贡献的个人和集体，均已在文中以明确方式标明。本人完全意识到本声明的法律结果由本人承担。

\vspace{2cm}
\begin{flushright}
    学位论文作者签名：\underline{\hspace{4cm}} \quad 日期：\underline{\hspace{3cm}}
\end{flushright}

\vspace{2cm}
\begin{center}
    {\zihao{3}\bfseries 学位论文版权使用授权书}
\end{center}

\vspace{1cm}
\zihao{-4}
本人完全了解学校关于学位论文知识产权的相关规定，同意学校保留并向国家有关部门或机构送交论文的复印件和电子版，允许论文被查阅和借阅。本人授权学校可将本学位论文的全部或部分内容编入有关数据库进行检索，可以采用影印、缩印或扫描等复制手段保存和汇编本学位论文。

\vspace{2cm}
\begin{flushright}
    学位论文作者签名：\underline{\hspace{4cm}} \quad 日期：\underline{\hspace{3cm}} \\[1cm]
    导师签名：\underline{\hspace{4.6cm}} \quad 日期：\underline{\hspace{3cm}}
\end{flushright}

% ==================== 摘要 ====================
\newpage
\pagenumbering{Roman}
\setcounter{page}{1}
\chapter*{摘\quad 要}
\addcontentsline{toc}{chapter}{摘要}

\zihao{-4}
这里是中文摘要。请简要介绍研究背景、研究目的、主要方法、核心结果和结论。摘要通常控制在500--800字左右。

本文的主要贡献包括：
\begin{enumerate}[label=(\arabic*)]
    \item 贡献一；
    \item 贡献二；
    \item 贡献三。
\end{enumerate}

\vspace{1cm}
\noindent\textbf{关键词：}关键词一；关键词二；关键词三

% ==================== Abstract ====================
\newpage
\chapter*{Abstract}
\addcontentsline{toc}{chapter}{Abstract}

\zihao{-4}
This is the English abstract. Please briefly introduce the research background, objectives, methods, main results, and conclusions. The abstract is usually around 300 words.

The main contributions of this thesis include:
\begin{enumerate}[label=(\arabic*)]
    \item Contribution one;
    \item Contribution two;
    \item Contribution three.
\end{enumerate}

\vspace{1cm}
\noindent\textbf{Keywords:} keyword one; keyword two; keyword three

% ==================== 目录 ====================
\newpage
\tableofcontents

% ==================== 插图索引（可选） ====================
\newpage
\listoffigures
\addcontentsline{toc}{chapter}{插图索引}

% ==================== 表格索引（可选） ====================
\newpage
\listoftables
\addcontentsline{toc}{chapter}{表格索引}

% ==================== 正文 ====================
\newpage
\cleardoublepage
\hypersetup{pageanchor=true}
\pagenumbering{arabic}
\setcounter{page}{1}

\chapter{绪论}
\label{chap:intro}

\section{研究背景与意义}
\label{sec:background}

请在此介绍研究背景，说明本课题的来源、研究意义以及研究的必要性。

\section{国内外研究现状}
\label{sec:related}

介绍与本研究相关的国内外研究进展，可引用相关文献\cite{vaswani2017attention,he2016deep}。

\section{研究内容与创新点}
\label{sec:content}

说明本文的主要研究内容、技术路线以及主要创新点。

\section{论文组织结构}
\label{sec:structure}

本文共分为X章，各章内容安排如下：

第\ref{chap:related}章介绍相关理论基础；

第\ref{chap:method}章阐述本文提出的方法；

第\ref{chap:experiment}章为实验设计与结果分析；

第\ref{chap:conclusion}章总结全文并展望未来工作。


\chapter{相关理论与技术基础}
\label{chap:related}

\section{基本概念}
\label{sec:concept}

介绍研究中涉及的基本概念和术语。

\section{相关技术}
\label{sec:technology}

介绍研究所依赖的核心技术、模型或方法。

\section{本章小结}
\label{sec:related-summary}

总结本章内容，引出下一章的研究方法。


\chapter{研究方法与设计}
\label{chap:method}

\section{问题定义}
\label{sec:problem}

明确本文要解决的问题，给出形式化定义。

\section{方法框架}
\label{sec:framework}

介绍所提方法的整体框架，可配合图\ref{fig:framework}说明。

\begin{figure}[htbp]
    \centering
    \includegraphics[width=0.8\textwidth]{example-image}
    \caption{方法整体框架}
    \label{fig:framework}
\end{figure}

\section{核心算法}
\label{sec:algorithm}

可以用算法伪代码描述核心流程，如算法\ref{alg:core}所示。

\begin{algorithm}[htbp]
    \caption{核心算法}
    \label{alg:core}
    \begin{algorithmic}[1]
        \Require 输入数据 $X$，参数 $\theta$
        \Ensure 输出结果 $Y$
        \State 初始化模型参数
        \For{每个训练轮次 $t$}
            \State 前向传播计算预测值
            \State 计算损失函数 $\mathcal{L}$
            \State 反向传播更新参数
        \EndFor
        \State \Return $Y$
    \end{algorithmic}
\end{algorithm}

\section{本章小结}
\label{sec:method-summary}

总结本章提出的方法及其优势。


\chapter{实验与结果分析}
\label{chap:experiment}

\section{实验设置}
\label{sec:setup}

介绍实验环境、数据集、评价指标和对比方法。

\section{实验结果}
\label{sec:result}

展示实验结果，可使用表格\ref{tab:result}进行对比。

\begin{table}[htbp]
    \centering
    \caption{实验结果对比}
    \label{tab:result}
    \begin{tabular}{lccc}
        \toprule
        方法 & 指标1 & 指标2 & 指标3 \\
        \midrule
        方法A & 0.85 & 0.82 & 0.80 \\
        方法B & 0.87 & 0.84 & 0.81 \\
        本文方法 & \textbf{0.91} & \textbf{0.88} & \textbf{0.86} \\
        \bottomrule
    \end{tabular}
\end{table}

\section{结果分析}
\label{sec:analysis}

对实验结果进行分析，讨论方法的有效性和优越性。

\section{本章小结}
\label{sec:experiment-summary}

总结实验结果。


\chapter{总结与展望}
\label{chap:conclusion}

\section{工作总结}
\label{sec:summary}

总结全文的主要研究内容和成果。

\section{未来展望}
\label{sec:future}

指出研究的不足之处，并对未来工作进行展望。


% ==================== 参考文献 ====================
\bibliography{references/references}

% ==================== 致谢 ====================
\chapter*{致\quad 谢}
\addcontentsline{toc}{chapter}{致谢}

\zihao{-4}
感谢导师的悉心指导，感谢家人和朋友的支持与鼓励，感谢评审老师的宝贵意见。

% ==================== 攻读学位期间取得的成果（可选） ====================
\chapter*{攻读硕士学位期间取得的成果}
\addcontentsline{toc}{chapter}{攻读硕士学位期间取得的成果}

\zihao{-4}
\begin{enumerate}[label={[\arabic*]}]
    \item 已发表论文或专利；
    \item 参与的研究项目；
    \item 获得的荣誉或奖励。
\end{enumerate}

\end{document}
